\section{Implementation Walkthrough}
\label{sec:impl}

This section reads the most load-bearing pieces of source. The full
codebase is available in the project repository; the listings below are
the parts a reader needs to internalise.

\subsection{Express bootstrap}

The API entry point composes middleware in a deliberate order: security
headers first (Helmet), then CORS, then the JSON body parser with a tight
size limit, then the structured logger, then the global rate limiter,
then the routers.

\begin{lstlisting}[caption={\texttt{server/server.js} (key fragment)},label={lst:server}]
app.set('trust proxy', 1);
app.use(helmet());
app.use(cors({ origin: process.env.CORS_ORIGIN || '*' }));
app.use(express.json({ limit: '256kb' }));
app.use(httpLogger);
app.use('/api', apiLimiter);

app.use('/api/auth',     authRoutes);
app.use('/api/projects', projectRoutes);
app.use('/api/grades',   gradeRoutes);
app.use('/api/votes',    voteRoutes);
app.use('/api/session',  sessionRoutes);
app.use('/api/schedule', scheduleRoutes);
app.use('/api/winners',  winnersRoutes);
app.use('/api/ai',       aiRoutes);
app.use('/api/admin',    adminRoutes);
\end{lstlisting}

\texttt{trust proxy} is set so that \texttt{req.ip} reflects the
\texttt{X-Forwarded-For} header set by Render's edge, not the loopback
address of Render's reverse proxy. Without this line, the per-IP rate
limit would mistakenly limit \emph{everyone} to one shared bucket.

\subsection{Validation middleware}

\begin{lstlisting}[caption={\texttt{server/middleware/validate.js}; the generic validator},label={lst:validate}]
export const validate = (schemas) => (req, res, next) => {
  try {
    if (schemas.body)   req.body   = schemas.body.parse(req.body ?? {});
    if (schemas.query)  req.query  = schemas.query.parse(req.query ?? {});
    if (schemas.params) req.params = schemas.params.parse(req.params ?? {});
    next();
  } catch (err) {
    if (err instanceof z.ZodError) {
      return res.status(400).json({
        message: 'Invalid request',
        issues: err.issues.map((i) => ({ path: i.path.join('.'), message: i.message })),
      });
    }
    next(err);
  }
};
\end{lstlisting}

The schemas object lives in the same file:

\begin{lstlisting}[caption={Schemas (excerpt)},label={lst:schemas}]
export const schemas = {
  login: {
    body: z.object({
      email: z.string().email().max(120),
      password: z.string().min(1).max(200),
    }),
  },
  voteCast: {
    body: z.object({
      projectId: z.string().uuid(),
      browserToken: z.string().uuid(),
      voteToken: z.string().min(10).max(400).optional(),
    }),
  },
  gradeBody: {
    body: z.object({
      scores: z.object({
        c1: z.number().min(0).max(5).nullable().optional(),
        c2: z.number().min(0).max(5).nullable().optional(),
        c3: z.number().min(0).max(5).nullable().optional(),
        c4: z.number().min(0).max(5).nullable().optional(),
        c5: z.number().min(0).max(5).nullable().optional(),
      }).default({}),
      feedback: z.string().max(5000).optional(),
    }),
    params: z.object({ projectId: z.string().uuid() }),
  },
  /* ... */
};
\end{lstlisting}

\subsection{Per-route rate limits}

\begin{lstlisting}[caption={\texttt{server/middleware/limits.js}},label={lst:limits}]
export const loginLimiter = rateLimit({
  windowMs: 60 * 1000, max: 5,
  standardHeaders: true, legacyHeaders: false,
  message: { message: 'Too many login attempts. Please wait a minute.' },
});

export const voteLimiter = rateLimit({
  windowMs: 60 * 1000, max: 5,
  standardHeaders: true, legacyHeaders: false,
  message: { message: 'Too many votes from this address.' },
});

export const gradeLimiter = rateLimit({
  windowMs: 60 * 1000, max: 60,
  standardHeaders: true, legacyHeaders: false,
  keyGenerator: (req) => req.judge?._id || req.ip,
  message: { message: 'Too many grading requests, slow down.' },
});
\end{lstlisting}

The grade limiter keys on \texttt{req.judge._id} rather than \texttt{req.ip}
so that two judges behind one campus NAT do not penalise each other.

\subsection{Cache helper}

\begin{lstlisting}[caption={\texttt{server/utils/cache.js}},label={lst:cache}]
import { LRUCache } from 'lru-cache';

const store = new LRUCache({ max: 1000, ttl: 60_000 });
let hits = 0, misses = 0;

export const cache = {
  async getOrSet(key, ttlMs, loader) {
    const cached = store.get(key);
    if (cached !== undefined) { hits++; return cached; }
    misses++;
    const fresh = await loader();
    store.set(key, fresh, { ttl: ttlMs });
    return fresh;
  },
  invalidate(prefix) {
    for (const k of store.keys()) {
      if (typeof k === 'string' && k.startsWith(prefix)) store.delete(k);
    }
  },
  stats() {
    return { size: store.size, hits, misses, hitRate: hits / (hits + misses || 1) };
  },
};
\end{lstlisting}

\subsection{HMAC vote token}

\begin{lstlisting}[caption={\texttt{server/utils/voteToken.js} (issue + verify)},label={lst:vote-token}]
export function issueVoteToken(browserToken) {
  const issuedAt = Date.now();
  const nonce = randomBytes(8).toString('hex');
  const payload = `${browserToken}|${issuedAt}|${nonce}`;
  const sig = createHmac('sha256', SECRET).update(payload).digest('hex');
  return `${Buffer.from(payload).toString('base64url')}.${sig}`;
}

export function verifyVoteToken(token, expectedBrowserToken) {
  if (typeof token !== 'string' || !token.includes('.')) return { ok: false, reason: 'malformed' };
  const [b64, sig] = token.split('.');
  let payload;
  try { payload = Buffer.from(b64, 'base64url').toString('utf8'); }
  catch { return { ok: false, reason: 'malformed' }; }
  const expectedSig = createHmac('sha256', SECRET).update(payload).digest('hex');
  if (sig.length !== expectedSig.length) return { ok: false, reason: 'bad-signature' };
  if (!timingSafeEqual(Buffer.from(sig), Buffer.from(expectedSig))) {
    return { ok: false, reason: 'bad-signature' };
  }
  const [browserToken, issuedAtStr] = payload.split('|');
  if (browserToken !== expectedBrowserToken) return { ok: false, reason: 'token-mismatch' };
  const issuedAt = Number(issuedAtStr);
  if (!Number.isFinite(issuedAt)) return { ok: false, reason: 'malformed' };
  if (Date.now() - issuedAt > TTL_MS) return { ok: false, reason: 'expired' };
  return { ok: true };
}
\end{lstlisting}

The use of \texttt{timingSafeEqual} closes the timing-attack vector;
without it, a sufficiently determined attacker could iteratively guess
signature bytes by measuring response times.

\subsection{Anomaly scoring}

\begin{lstlisting}[caption={\texttt{server/services/anomalyDetector.js} (scoring)},label={lst:anom}]
export function scoreVote({ ip, userAgent, browserToken }) {
  trim();
  const reasons = [];
  let score = 0;
  const now = Date.now();
  const subnet = subnet24(ip);

  const sameIp = recent.filter((r) => r.ip === ip);
  const sameSubnet = recent.filter((r) => subnet24(r.ip) === subnet);
  const sameUa = recent.filter((r) => r.ua === userAgent);
  const burstIp = sameIp.filter((r) => now - r.time < BURST_WINDOW_MS);
  const burstSubnet = sameSubnet.filter((r) => now - r.time < BURST_WINDOW_MS);

  if (burstIp.length >= 3)        { score += 40; reasons.push(`burst:${burstIp.length} from same IP within 60s`); }
  else if (burstIp.length >= 1)   { score += 15; reasons.push(`recent vote from same IP within 60s`); }
  if (burstSubnet.length >= 8)    { score += 30; reasons.push(`burst:${burstSubnet.length} from /24 within 60s`); }
  if (sameUa.length >= 6)         { score += 15; reasons.push(`UA reused ${sameUa.length}x in 5min`); }
  if (sameIp.length >= 12)        { score += 20; reasons.push(`${sameIp.length} votes from same IP in 5min`); }
  if (!userAgent || userAgent.length < 10) { score += 10; reasons.push('thin/missing User-Agent'); }

  let severity = 'none';
  if (score >= 60) severity = 'high';
  else if (score >= 30) severity = 'medium';
  else if (score >= 10) severity = 'low';

  recent.push({ ip, ua: userAgent, time: now, browserToken });
  return { score, severity, reasons };
}
\end{lstlisting}

\subsection{Vote controller (the cold path)}

\begin{lstlisting}[caption={\texttt{server/controllers/voteController.js} (castVote, abridged)},label={lst:cast}]
export const castVote = async (req, res) => {
  const { projectId, browserToken, voteToken } = req.body;
  const ip = req.ip ?? '';
  const userAgent = req.headers['user-agent'] || '';

  if (process.env.REQUIRE_VOTE_TOKEN === 'true') {
    const v = verifyVoteToken(voteToken, browserToken);
    if (!v.ok) return res.status(403).json({ message: `Vote token invalid: ${v.reason}` });
  }

  const { data: project } = await supabase.from('projects').select('id').eq('id', projectId).single();
  if (!project) return res.status(404).json({ message: 'Project not found' });

  const { data: existing } = await supabase
    .from('votes').select('project_id').eq('browser_token', browserToken).maybeSingle();
  if (existing) {
    return res.status(409).json({ message: 'Already voted', votedProjectId: existing.project_id });
  }

  const anomaly = scoreVote({ ip, userAgent, browserToken });

  const { data: insertedVote, error: insertErr } = await supabase
    .from('votes').insert({ project_id: projectId, browser_token: browserToken, ip_address: ip })
    .select().single();
  if (insertErr) return res.status(500).json({ message: 'Could not record vote' });

  await supabase.rpc('increment_vote_count', { p_id: projectId });

  if (anomaly.severity !== 'none') {
    recordFlag({ voteId: insertedVote.id, projectId, browserToken, ip,
                 userAgent: hashUserAgent(userAgent),
                 score: anomaly.score, severity: anomaly.severity, reasons: anomaly.reasons })
      .catch((e) => logger.warn({ err: e.message }, 'flag persist failed'));
  }

  cache.invalidate('projects:'); cache.invalidate('winners:'); cache.invalidate('vote-counts');
  broadcast('vote.cast', { projectId, flagged: anomaly.severity !== 'none', severity: anomaly.severity });

  res.status(201).json({ message: 'Vote recorded', projectId, flagged: anomaly.severity !== 'none' });
};
\end{lstlisting}

The hot path is exactly what the architecture promised: one
\texttt{SELECT} for project existence, one \texttt{SELECT} for duplicate
detection, one in-memory anomaly score, one \texttt{INSERT}, one RPC,
one cache invalidation, one broadcast. No N+1, no joins, no hidden
fan-out.

\subsection{SSE event bus}

\begin{lstlisting}[caption={\texttt{server/services/eventBus.js}},label={lst:bus}]
const clients = new Set();

export function addClient(res)    { clients.add(res); }
export function removeClient(res) { clients.delete(res); }

export function broadcast(event, data) {
  const payload = `event: ${event}\ndata: ${JSON.stringify(data)}\n\n`;
  for (const res of clients) {
    try { res.write(payload); }
    catch (err) { clients.delete(res); }
  }
}

setInterval(() => {
  for (const res of clients) {
    try { res.write(': ping\n\n'); } catch { clients.delete(res); }
  }
}, 25_000).unref?.();
\end{lstlisting}

The route handler that turns a normal HTTP request into an SSE stream:

\begin{lstlisting}[caption={SSE handler in \texttt{sessionController.js}},label={lst:sse}]
export const streamSession = (req, res) => {
  res.set({
    'Content-Type': 'text/event-stream',
    'Cache-Control': 'no-cache, no-transform',
    Connection: 'keep-alive',
    'X-Accel-Buffering': 'no',
  });
  res.flushHeaders?.();

  readMain()
    .then((row) => res.write(`event: session.snapshot\ndata: ${JSON.stringify(shapeSession(row))}\n\n`))
    .catch((err) => logger.warn({ err: err.message }, 'sse snapshot failed'));

  addClient(res);
  req.on('close', () => removeClient(res));
};
\end{lstlisting}

\subsection{Claude API call with prompt caching}

\begin{lstlisting}[caption={\texttt{server/services/claude.js} (the message create)},label={lst:claude}]
const RUBRIC_SYSTEM = `You are an experienced judge for a university soft-skills competition.
... (rubric block) ...`;

export async function generateInsights(project) {
  if (!ENABLE_REAL) return stubInsights(project);

  const userText = [
    `Project: ${project.title}`,
    `Segment: ${project.segmentType}`,
    `Team / speaker: ${project.teamName || '(unknown)'}`,
    `Category: ${project.category || '(unknown)'}`,
    `Tags: ${(project.tags || []).join(', ') || '(none)'}`,
    '',
    'Description:',
    project.description || '(no description provided)',
  ].join('\n');

  const resp = await client.messages.create({
    model: MODEL,
    max_tokens: 700,
    system: [
      { type: 'text', text: RUBRIC_SYSTEM, cache_control: { type: 'ephemeral' } },
    ],
    messages: [{ role: 'user', content: userText }],
  });

  const text = resp.content.filter((b) => b.type === 'text').map((b) => b.text).join('\n').trim();
  logger.info({
    cacheRead:  resp.usage?.cache_read_input_tokens ?? 0,
    cacheWrite: resp.usage?.cache_creation_input_tokens ?? 0,
    input:  resp.usage?.input_tokens ?? 0,
    output: resp.usage?.output_tokens ?? 0,
  }, 'claude.insights.usage');
  return extractJson(text);
}
\end{lstlisting}

The two log fields \texttt{cacheRead} and \texttt{cacheWrite} let us
verify on every call that the prompt cache is doing its job; in steady
state, \texttt{cacheRead} is large and \texttt{cacheWrite} is zero.

\subsection{Audit helper}

\begin{lstlisting}[caption={\texttt{server/utils/audit.js}},label={lst:audit}]
export async function audit(entry) {
  try {
    const { error } = await supabase.from('audit_logs').insert({
      actor_type: entry.actorType,
      actor_id:   entry.actorId ?? null,
      action:     entry.action,
      target_type: entry.targetType ?? null,
      target_id:   entry.targetId ?? null,
      metadata:    entry.metadata ?? {},
      ip:          entry.ip ?? null,
    });
    if (error) logger.warn({ err: error.message, action: entry.action }, 'audit insert failed');
  } catch (e) {
    logger.warn({ err: e.message, action: entry.action }, 'audit threw');
  }
}
\end{lstlisting}

The audit insert is best-effort: a failure to persist the audit row
should not block the user-facing operation. We log at WARN so an
operator scanning logs can still see the gap.

\subsection{Client SSE hook}

\begin{lstlisting}[caption={\texttt{client/src/hooks/useEventStream.js} (key fragment)},label={lst:hook}]
useEffect(() => {
  if (!url) return undefined;
  let es = null, watchdog = null, lastEventAt = Date.now(), stopped = false;
  const staleMs = opts.staleMs ?? 45_000;

  function connect() {
    if (stopped) return;
    try { es?.close(); } catch { /* noop */ }
    es = new EventSource(url);
    es.onopen  = () => { lastEventAt = Date.now(); };
    es.onerror = () => { /* native retry handles transient errors */ };
    es.onmessage = (ev) => {
      lastEventAt = Date.now();
      try { handlersRef.current?.message?.(JSON.parse(ev.data)); } catch { /* ignore */ }
    };
    for (const name of Object.keys(handlersRef.current || {})) {
      if (name === 'message') continue;
      es.addEventListener(name, (ev) => {
        lastEventAt = Date.now();
        try { handlersRef.current[name]?.(ev.data ? JSON.parse(ev.data) : null, ev); }
        catch { handlersRef.current[name]?.(null, ev); }
      });
    }
  }

  connect();
  watchdog = setInterval(() => {
    if (Date.now() - lastEventAt > staleMs) { lastEventAt = Date.now(); connect(); }
  }, Math.max(staleMs / 2, 5_000));

  return () => { stopped = true; clearInterval(watchdog); try { es?.close(); } catch {} };
}, [url, opts.staleMs]);
\end{lstlisting}

\subsection{Frontend component map}

The page-level component graph is shown in
Figure~\ref{fig:class-frontend}. Two providers wrap the router:
\texttt{AuthProvider} owns the JWT and judge profile; \texttt{LiveSessionProvider}
owns the SSE subscription and the session snapshot.

\fig{class-frontend}{0.95}{Top-level frontend component hierarchy.}
